% !TEX program = xelatex
% LingTu (灵途) 产品手册 v2 — 面向客户，产品语言，无技术术语
% 编译: xelatex lingtu_product_manual.tex (需要两次以生成目录)

\documentclass[a4paper, 12pt, openany]{book}

% ══════════════════════════════════════════════════════════════
% 基础包
% ══════════════════════════════════════════════════════════════
\usepackage[UTF8, heading=true]{ctex}
\usepackage[margin=2.5cm, top=3cm, bottom=3cm, headheight=14pt]{geometry}
\usepackage{graphicx}
\usepackage{xcolor}
\usepackage{tikz}
\usetikzlibrary{shapes.geometric, arrows.meta, positioning, calc, fit, backgrounds, decorations.pathreplacing, shadows.blur, patterns}
\usepackage{booktabs}
\usepackage{longtable}
\usepackage{tabularx}
\usepackage{enumitem}
\usepackage{fancyhdr}
\usepackage{titlesec}
\usepackage{tcolorbox}
\tcbuselibrary{skins, breakable}
\usepackage{fontspec}
\usepackage{hyperref}
\usepackage{pifont}
\usepackage{multicol}
\usepackage{float}
\usepackage[section]{placeins}  % 浮动体不跨section
\usepackage{caption}
\usepackage{subcaption}
\usepackage{amssymb}
\usepackage{amsmath}
\usepackage{setspace}
\usepackage{wrapfig}

% 排版优化：紧凑排版
\raggedbottom
\setlength{\intextsep}{4pt plus 1pt minus 1pt}
\setlength{\floatsep}{4pt plus 1pt minus 1pt}
\setlength{\textfloatsep}{6pt plus 2pt minus 2pt}
\setlength{\parskip}{2pt plus 1pt}
\setlength{\abovecaptionskip}{4pt}
\setlength{\belowcaptionskip}{2pt}
\setlength{\topsep}{2pt}
\setlength{\partopsep}{0pt}
\setlength{\itemsep}{1pt}

% ══════════════════════════════════════════════════════════════
% 颜色定义
% ══════════════════════════════════════════════════════════════
\definecolor{qpblue}{RGB}{0, 82, 155}
\definecolor{qpdark}{RGB}{30, 40, 55}
\definecolor{qplight}{RGB}{230, 240, 250}
\definecolor{qpaccent}{RGB}{0, 150, 136}
\definecolor{qpwarn}{RGB}{255, 152, 0}
\definecolor{qpgray}{RGB}{120, 130, 140}
\definecolor{qpgreen}{RGB}{46, 125, 50}
\definecolor{qporange}{RGB}{230, 126, 34}
\definecolor{qpred}{RGB}{198, 40, 40}
\definecolor{scenebg}{RGB}{245, 248, 252}
\definecolor{cardblue}{RGB}{227, 242, 253}
\definecolor{cardgreen}{RGB}{232, 245, 233}
\definecolor{cardorange}{RGB}{255, 243, 224}
\definecolor{cardpurple}{RGB}{237, 231, 246}
\definecolor{cardred}{RGB}{255, 235, 238}
\definecolor{cardteal}{RGB}{224, 242, 241}

% ══════════════════════════════════════════════════════════════
% 超链接与排版
% ══════════════════════════════════════════════════════════════
\hypersetup{
  colorlinks=true,
  linkcolor=qpblue,
  urlcolor=qpaccent,
  bookmarksdepth=3,
  pdfauthor={上海穹沛科技有限公司},
  pdftitle={灵途 (LingTu) 自主导航系统 产品手册},
}

% ══════════════════════════════════════════════════════════════
% 页眉页脚
% ══════════════════════════════════════════════════════════════
\pagestyle{fancy}
\fancyhf{}
\fancyhead[L]{\small\color{qpgray}\leftmark}
\fancyhead[R]{\small\color{qpgray}灵途自主导航系统}
\fancyfoot[C]{\small\color{qpgray}\thepage}
\renewcommand{\headrulewidth}{0.4pt}
\renewcommand{\footrulewidth}{0pt}

% ══════════════════════════════════════════════════════════════
% 章节标题样式
% ══════════════════════════════════════════════════════════════
\ctexset{
  chapter = {
    format = \huge\bfseries\color{qpblue},
    nameformat = {},
    titleformat = {},
    beforeskip = 8pt,
    afterskip = 8pt,
    break = {},
  },
  section = {
    format = \Large\bfseries\color{qpdark},
    beforeskip = 8pt,
    afterskip = 4pt,
  },
  subsection = {
    format = \large\bfseries\color{qpdark},
    beforeskip = 6pt,
    afterskip = 3pt,
  },
}

% ══════════════════════════════════════════════════════════════
% 自定义框
% ══════════════════════════════════════════════════════════════
\newtcolorbox{infobox}[1][]{
  colback=qplight, colframe=qpblue, coltitle=white, fonttitle=\bfseries,
  title=#1, rounded corners, boxrule=0.8pt,
  left=6pt, right=6pt, top=4pt, bottom=4pt, breakable,
  before skip=4pt, after skip=4pt,
}

\newtcolorbox{tipbox}[1][]{
  colback=green!5, colframe=qpaccent, coltitle=white, fonttitle=\bfseries,
  title=#1, rounded corners, boxrule=0.8pt,
  left=6pt, right=6pt, top=4pt, bottom=4pt, breakable,
  before skip=4pt, after skip=4pt,
}

\newtcolorbox{warnbox}[1][]{
  colback=orange!5, colframe=qpwarn, coltitle=white, fonttitle=\bfseries,
  title=#1, rounded corners, boxrule=0.8pt,
  left=6pt, right=6pt, top=4pt, bottom=4pt, breakable,
  before skip=4pt, after skip=4pt,
}

\newtcolorbox{scenebox}[1][]{
  colback=scenebg, colframe=qpgray!40, fonttitle=\bfseries\color{qpdark},
  title=#1, rounded corners, boxrule=0.5pt,
  left=6pt, right=6pt, top=4pt, bottom=4pt, breakable,
  before skip=4pt, after skip=4pt,
}

% 功能卡片
\newcommand{\featurecard}[4]{%
  \begin{tcolorbox}[
    colback=#1, colframe=#1!60!black!30, boxrule=0.5pt,
    rounded corners, left=6pt, right=6pt, top=4pt, bottom=4pt,
    width=\linewidth,
  ]
  {\bfseries\color{qpdark}#2} \hfill {\small\color{qpgray}#3} \\[3pt]
  {\small #4}
  \end{tcolorbox}
}

% ══════════════════════════════════════════════════════════════
% 文档开始
% ══════════════════════════════════════════════════════════════
\begin{document}

% ──────────────────────────────────────────────────────────────
% 封面
% ──────────────────────────────────────────────────────────────
\begin{titlepage}
\begin{tikzpicture}[remember picture, overlay]
  \fill[qpdark] (current page.south west) rectangle (current page.north east);
  \fill[qpblue] ([yshift=-4cm]current page.north west) rectangle ([yshift=-4.15cm]current page.north east);
  \fill[qpaccent] ([yshift=-4.15cm]current page.north west) rectangle ([yshift=-4.25cm]current page.north east);
  \fill[qpblue, opacity=0.3] ([yshift=3cm]current page.south west) rectangle ([yshift=3.1cm]current page.south east);
\end{tikzpicture}

\vspace*{3cm}
\begin{center}
  {\fontsize{18}{22}\selectfont\color{qpaccent}\textbf{上海穹沛科技有限公司}}

  \vspace{1.5cm}

  {\fontsize{48}{56}\selectfont\color{white}\textbf{灵途}}

  \vspace{0.5cm}

  {\fontsize{20}{24}\selectfont\color{white!80}LingTu Autonomous Navigation System}

  \vspace{0.8cm}

  {\fontsize{16}{20}\selectfont\color{qpaccent}产\quad 品\quad 手\quad 册}

  \vspace{3cm}

  {\large\color{white!60}四足机器人自主导航解决方案}

  \vspace{0.3cm}

  {\large\color{white!60}语义理解 \quad $\cdot$ \quad 自主决策 \quad $\cdot$ \quad 安全可靠}

  \vfill

  {\color{white!40}\rule{8cm}{0.5pt}}

  \vspace{0.5cm}

  {\normalsize\color{white!50}版本 1.8.0 \quad | \quad 2026 年 3 月}

  \vspace{1cm}

  {\small\color{white!30}本文档包含穹沛科技机密信息，仅限授权人员阅读}
\end{center}
\end{titlepage}

% ──────────────────────────────────────────────────────────────
% 版本页
% ──────────────────────────────────────────────────────────────
\thispagestyle{empty}
\vspace*{\fill}
\begin{center}
\begin{tabular}{ll}
  \toprule
  \textbf{文档信息} & \\
  \midrule
  文档名称 & 灵途自主导航系统 产品手册 \\
  版本号 & v1.8.0 \\
  发布日期 & 2026 年 3 月 15 日 \\
  编制单位 & 上海穹沛科技有限公司 \\
  \bottomrule
\end{tabular}
\end{center}
\vspace*{\fill}
\clearpage

% ──────────────────────────────────────────────────────────────
% 目录
% ──────────────────────────────────────────────────────────────
\tableofcontents
\clearpage


% ══════════════════════════════════════════════════════════════
%  第一章：产品概述
% ══════════════════════════════════════════════════════════════
\chapter{产品概述}

\section{灵途是什么}

\textbf{灵途}（LingTu）是由上海穹沛科技有限公司自主研发的\textbf{四足机器人自主导航系统}。

传统四足机器人就像一辆没有自动驾驶的汽车——功能强大，但每一步都需要人来操控。灵途做的事情，就是给四足机器人装上"自动驾驶大脑"：让它能自己看路、自己找目标、自己规划路线、自己避开障碍，最终自主完成任务。

更重要的是，灵途能听懂人话。操作员不需要在地图上精确点击坐标，只要对机器人说"去大门那边看看"或者"帮我找一下灭火器"，它就能理解指令、自动执行。

\begin{infobox}[一句话总结]
灵途让四足机器人从"需要全程遥控的工具"变成"说一句话就能自己干活的助手"。一个操作员可以同时管理多台机器人，大幅降低人力成本。
\end{infobox}

\section{解决什么问题}

在工厂巡检、园区安防、户外探测等场景中，四足机器人相比轮式机器人有着天然的地形适应优势——它能翻越台阶、穿过草地、应对复杂地形。但传统的四足机器人面临一个核心问题：\textbf{需要专人全程遥控操作}。

\begin{figure}[!htb]
\centering
\begin{tikzpicture}[
  pain/.style={draw=qpred!60, fill=cardred, rounded corners=6pt, minimum width=4.2cm, minimum height=2.2cm, align=center, font=\small},
  solve/.style={draw=qpgreen!60, fill=cardgreen, rounded corners=6pt, minimum width=4.2cm, minimum height=2.2cm, align=center, font=\small},
  arr/.style={-{Stealth[length=6pt]}, thick, qpblue},
]
  % 痛点
  \node[pain] (p1) at (0, 0) {\textbf{人力成本高}\\[2pt]每台机器人\\需要一名操作员\\全程关注};
  \node[pain] (p2) at (5, 0) {\textbf{操作门槛高}\\[2pt]遥控四足机器人\\需要专业培训\\难以普及};
  \node[pain] (p3) at (10, 0) {\textbf{效率低下}\\[2pt]人工遥控速度慢\\容易遗漏区域\\无法24小时工作};

  % 箭头
  \draw[arr] (0, -1.5) -- (0, -2.5);
  \draw[arr] (5, -1.5) -- (5, -2.5);
  \draw[arr] (10, -1.5) -- (10, -2.5);

  % 灵途标签
  \node[fill=qpblue, text=white, rounded corners=3pt, font=\small\bfseries, minimum width=12.5cm] at (5, -2) {安装灵途后};

  % 解决方案
  \node[solve] (s1) at (0, -4) {\textbf{一人管多机}\\[2pt]机器人自主执行\\操作员只需监控\\人力降低80\%};
  \node[solve] (s2) at (5, -4) {\textbf{说话就能用}\\[2pt]自然语言指令\\普通员工即可\\操作零门槛};
  \node[solve] (s3) at (10, -4) {\textbf{全天候运行}\\[2pt]自主巡检不停歇\\自动避障不遗漏\\效率提升5倍};
\end{tikzpicture}
\caption{灵途解决的核心痛点}
\end{figure}

\section{适用场景}

灵途适用于需要四足机器人在户外或半户外环境中自主移动的各种场景。以下是典型应用：

\begin{figure}[!htb]
\centering
\begin{tikzpicture}[
  scene/.style={draw=qpgray!20, fill=white, rounded corners=6pt, minimum width=6.8cm, minimum height=2.6cm, align=left, font=\small, inner sep=8pt},
]
  % 第一行
  \node[scene] (s1) at (0, 0) {
    {\large\bfseries\color{qpdark} 工厂巡检} \\[4pt]
    在车间、仓库、产线之间定时巡检。\\
    检查设备状态、安全隐患、消防设施。\\
    {\color{qpgray}\small 典型指令："去检查3号车间的灭火器"}
  };
  \node[scene] (s2) at (7.5, 0) {
    {\large\bfseries\color{qpdark} 园区安防} \\[4pt]
    办公园区、物流园区24小时安全巡逻。\\
    异常检测、入侵预警、夜间巡视。\\
    {\color{qpgray}\small 典型指令："沿围栏巡逻一圈"}
  };

  % 第二行
  \node[scene] (s3) at (0, -3.6) {
    {\large\bfseries\color{qpdark} 电力巡线} \\[4pt]
    沿电力线路巡视，检查塔基设备。\\
    复杂地形自动适应，定点停留拍照。\\
    {\color{qpgray}\small 典型指令："沿路线巡检到终点"}
  };
  \node[scene] (s4) at (7.5, -3.6) {
    {\large\bfseries\color{qpdark} 建筑工地} \\[4pt]
    施工现场安全检查和进度监控。\\
    动态环境自动避障，实时画面回传。\\
    {\color{qpgray}\small 典型指令："去查看北侧脚手架"}
  };

  % 第三行
  \node[scene] (s5) at (0, -7.2) {
    {\large\bfseries\color{qpdark} 应急救援} \\[4pt]
    灾害现场环境探测和搜救辅助。\\
    未知环境自主探索，实时回传信息。\\
    {\color{qpgray}\small 典型指令："探索前方区域"}
  };
  \node[scene] (s6) at (7.5, -7.2) {
    {\large\bfseries\color{qpdark} 科研教育} \\[4pt]
    机器人导航算法研究和教学实验。\\
    开放接口、可扩展、支持二次开发。\\
    {\color{qpgray}\small 提供完整的API和开发文档}
  };
\end{tikzpicture}
\caption{灵途适用场景一览}
\end{figure}


% ══════════════════════════════════════════════════════════════
%  第二章：产品功能全景
% ══════════════════════════════════════════════════════════════
\chapter{产品功能全景}

灵途系统由六大功能模块组成，从"看得见"到"走得到"到"管得住"，形成完整的自主导航闭环。

\begin{figure}[!htb]
\centering
\begin{tikzpicture}[
  mod/.style={draw=qpblue!40, fill=#1, rounded corners=10pt, minimum width=4cm, minimum height=3.5cm, align=center, font=\small, inner sep=6pt},
  conn/.style={-{Stealth[length=5pt]}, thick, qpblue!50},
]
  % 六大模块
  \node[mod=cardblue] (see) at (0, 4) {
    {\bfseries\large\color{qpdark} 智能感知}\\[6pt]
    AI视觉识别 125+ 类物体\\
    激光雷达3D环境扫描\\
    地形分析与障碍检测\\
    实时目标跟踪
  };
  \node[mod=cardgreen] (map) at (5.5, 4) {
    {\bfseries\large\color{qpdark} 环境建图}\\[6pt]
    一键三维地图构建\\
    厘米级定位精度\\
    抗振动建图技术\\
    地图保存与管理
  };
  \node[mod=cardorange] (nav) at (11, 4) {
    {\bfseries\large\color{qpdark} 自主导航}\\[6pt]
    全局路径自动规划\\
    实时动态避障\\
    卡死自动检测恢复\\
    多航点任务执行
  };
  \node[mod=cardpurple] (sem) at (0, -0.5) {
    {\bfseries\large\color{qpdark} 语义理解}\\[6pt]
    自然语言指令输入\\
    "快速+深度"双模决策\\
    场景图语义推理\\
    智能失败恢复
  };
  \node[mod=cardred] (safe) at (5.5, -0.5) {
    {\bfseries\large\color{qpdark} 安全防护}\\[6pt]
    8层纵深安全架构\\
    硬件级急停覆盖\\
    近场避障保护\\
    三环实时监控
  };
  \node[mod=cardteal] (ctrl) at (11, -0.5) {
    {\bfseries\large\color{qpdark} 远程管控}\\[6pt]
    手机/电脑App控制\\
    Web实时监控仪表盘\\
    远程OTA软件升级\\
    多机集中管理
  };

  % 连接
  \draw[conn] (see) -- (map);
  \draw[conn] (map) -- (nav);
  \draw[conn] (see) -- (sem);
  \draw[conn] (sem) -- (nav);
  \draw[conn] (nav) -- (safe);
  \draw[conn] (safe) -- (ctrl);
  \draw[conn] (see) -- (safe);
\end{tikzpicture}
\caption{灵途六大功能模块}
\end{figure}

\vspace{0.5cm}

以下章节将逐一详细介绍每个模块的功能和使用方式。


% ══════════════════════════════════════════════════════════════
%  第三章：智能建图
% ══════════════════════════════════════════════════════════════
\chapter{智能建图}

在机器人开始自主导航之前，需要先"认识"工作环境——就像新员工入职时先熟悉办公楼一样。灵途的建图功能让这个过程简单、快速、精确。

\section{建图流程}

\begin{figure}[!htb]
\centering
\begin{tikzpicture}[
  stp/.style={draw=qpblue!40, fill=cardblue, rounded corners=8pt, minimum width=2.8cm, minimum height=1.8cm, align=center, font=\small},
  num/.style={circle, fill=qpblue, text=white, font=\bfseries, minimum size=0.7cm},
  arr/.style={-{Stealth[length=6pt]}, thick, qpblue!60},
]
  \node[stp] (s1) at (0, 0) {打开App\\选择"建图模式"};
  \node[stp] (s2) at (3.8, 0) {遥控机器人\\走一圈工作区域};
  \node[stp] (s3) at (7.6, 0) {系统自动生成\\三维环境地图};
  \node[stp] (s4) at (11.4, 0) {一键保存\\后续直接使用};

  \node[num] at ([yshift=1.3cm]s1.north) {1};
  \node[num] at ([yshift=1.3cm]s2.north) {2};
  \node[num] at ([yshift=1.3cm]s3.north) {3};
  \node[num] at ([yshift=1.3cm]s4.north) {4};

  \draw[arr] (s1) -- (s2);
  \draw[arr] (s2) -- (s3);
  \draw[arr] (s3) -- (s4);

  \node[font=\footnotesize\color{qpgray}] at (1.9, -1.3) {用遥控器或App};
  \node[font=\footnotesize\color{qpgray}] at (5.7, -1.3) {全程自动处理};
  \node[font=\footnotesize\color{qpgray}] at (9.5, -1.3) {存储在机器人本地};
\end{tikzpicture}
\caption{建图只需四步}
\end{figure}

操作非常简单：打开手机App，选择"建图模式"，然后用遥控器或App摇杆操控机器人以正常步速在工作区域内走一圈。在行走过程中，机器人的激光雷达会自动扫描周围环境，系统实时将扫描数据拼合成一张完整的三维地图。走完一圈后，点击"保存地图"按钮即可。

\section{建图质量}

灵途的建图技术专门针对四足机器人进行了优化。四足机器人在行走时不可避免地产生身体晃动——这是轮式机器人不会遇到的挑战。灵途采用先进的"激光+惯性"融合定位技术，即使在机器人快速行走、身体剧烈晃动的情况下，也能保持极高的建图精度。

\begin{figure}[!htb]
\centering
\begin{tikzpicture}[
  metric/.style={draw=qpblue!30, fill=cardblue, rounded corners=6pt, minimum width=4.5cm, minimum height=2cm, align=center, font=\small},
]
  \node[metric] (m1) at (0, 0) {\textbf{\large 0.008 米}\\[4pt]标准建图精度\\（误差不到1厘米）};
  \node[metric] (m2) at (5.5, 0) {\textbf{\large 0.006 米}\\[4pt]增强抗振模式\\（极端振动场景）};
  \node[metric] (m3) at (11, 0) {\textbf{\large 15--20 分钟}\\[4pt]2000平方米\\典型建图时间};
\end{tikzpicture}
\caption{建图性能指标}
\end{figure}

系统提供两种定位算法供用户选择：\textbf{标准模式}适用于绝大多数场景；\textbf{增强抗振模式}在机器人高速奔跑或经过极度颠簸地形时提供更好的稳定性。切换方式非常简单，在启动时选择对应的配置即可。

\section{地图管理}

构建好的地图保存在机器人本地的高速存储中。用户可以通过App进行以下地图管理操作：

\begin{itemize}[leftmargin=2em, itemsep=4pt]
  \item \textbf{保存地图}：建图完成后一键保存，支持为地图命名（如"3号车间-2026年3月"）；
  \item \textbf{加载地图}：导航前选择要使用的地图，系统自动加载并定位；
  \item \textbf{多地图管理}：支持存储多张地图，适用于多个工作场地；
  \item \textbf{远程更新}：地图文件可以通过OTA远程下发到机器人。
\end{itemize}

\begin{tipbox}[什么时候需要重建地图？]
如果工作环境发生了较大的结构性变化（如搬迁了大型设备、修建了新墙），建议重新建图。对于日常的小变化（如移动了几把椅子），系统的实时避障功能可以自动应对，无需更新地图。
\end{tipbox}


% ══════════════════════════════════════════════════════════════
%  第四章：自主导航
% ══════════════════════════════════════════════════════════════
\chapter{自主导航}

建图完成后，机器人就可以在已知环境中自主导航了。用户只需告诉它目的地，剩下的全部交给灵途。

\section{导航方式}

灵途支持三种导航指令方式，满足不同使用习惯：

\begin{figure}[!htb]
\centering
\begin{tikzpicture}[
  way/.style={draw=qpblue!30, fill=#1, rounded corners=8pt, minimum width=4cm, minimum height=3cm, align=center, font=\small, inner sep=8pt},
]
  \node[way=cardblue] (w1) at (0, 0) {
    {\large\bfseries\color{qpblue} 地图点击}\\[6pt]
    在App地图上\\
    直接点击目标位置\\[4pt]
    {\color{qpgray}\footnotesize 最直观}
  };
  \node[way=cardgreen] (w2) at (5.5, 0) {
    {\large\bfseries\color{qpgreen} 语音/文字指令}\\[6pt]
    说出或输入目标描述\\
    如"去找灭火器"\\[4pt]
    {\color{qpgray}\footnotesize 最便捷}
  };
  \node[way=cardorange] (w3) at (11, 0) {
    {\large\bfseries\color{qporange} 巡检路线}\\[6pt]
    预设多个检查点\\
    机器人按顺序巡逻\\[4pt]
    {\color{qpgray}\footnotesize 最高效}
  };
\end{tikzpicture}
\caption{三种导航指令方式}
\end{figure}

\section{导航过程}

当机器人收到一个导航目标后，系统内部会自动完成一系列复杂的计算和决策，但对用户来说，整个过程是透明的——只需等待机器人到达目的地即可。

\begin{figure}[!htb]
\centering
\begin{tikzpicture}[
  phase/.style={draw=qpblue!40, fill=white, rounded corners=6pt, minimum width=12cm, minimum height=1.4cm, align=left, font=\small, inner xsep=10pt},
  icon/.style={circle, fill=#1, text=white, font=\bfseries\footnotesize, minimum size=0.9cm},
  timelab/.style={font=\footnotesize\color{qpaccent}, anchor=west},
]
  \node[phase] (p1) at (0, 5) {\hspace{1.2cm}\textbf{路线规划}：在地图上计算从当前位置到目标的最优路径};
  \node[icon=qpblue] at ([xshift=0.5cm]p1.west) {1};
  \node[timelab] at ([xshift=0.3cm]p1.east) {5--10 毫秒};

  \node[phase] (p2) at (0, 3.2) {\hspace{1.2cm}\textbf{路线分段}：将长路线分成多个中间航点，逐段推进};
  \node[icon=qpblue] at ([xshift=0.5cm]p2.west) {2};

  \node[phase] (p3) at (0, 1.4) {\hspace{1.2cm}\textbf{实时避障}：行走过程中持续扫描前方环境，动态躲避新出现的障碍};
  \node[icon=qpaccent] at ([xshift=0.5cm]p3.west) {3};
  \node[timelab] at ([xshift=0.3cm]p3.east) {持续运行};

  \node[phase] (p4) at (0, -0.4) {\hspace{1.2cm}\textbf{卡死恢复}：如果5秒内几乎没有前进，发出预警；10秒仍无进展，自动重新规划};
  \node[icon=qpwarn] at ([xshift=0.5cm]p4.west) {4};

  \node[phase] (p5) at (0, -2.2) {\hspace{1.2cm}\textbf{到达确认}：抵达目标位置后，系统自动报告"已到达"};
  \node[icon=qpgreen] at ([xshift=0.5cm]p5.west) {5};

  \draw[-{Stealth[length=4pt]}, gray!40, thick] (p1.south) -- (p2.north);
  \draw[-{Stealth[length=4pt]}, gray!40, thick] (p2.south) -- (p3.north);
  \draw[-{Stealth[length=4pt]}, gray!40, thick] (p3.south) -- (p4.north);
  \draw[-{Stealth[length=4pt]}, gray!40, thick] (p4.south) -- (p5.north);
\end{tikzpicture}
\caption{自主导航五阶段流程}
\end{figure}

灵途的路线规划速度极快——从收到目标到计算出完整路线，仅需 5--10 毫秒。规划出的路线会自动避开地图中已知的障碍区域，并优先选择平坦、安全的地面。

在行走过程中，系统不会"闭着眼睛走路"。它会持续分析前方的地形数据，如果发现有新出现的障碍物（比如一个临时放在路上的箱子），会自动调整路线绕行，整个过程无需人工干预。

\section{巡检任务}

对于需要定期执行的巡检工作，灵途支持\textbf{多航点巡检任务}：

\begin{enumerate}[leftmargin=2em, itemsep=4pt]
  \item 在App中依次点击巡检路线上的各个检查点；
  \item 设置巡检参数（如每个点停留时间、巡检频率）；
  \item 一键启动，机器人按顺序访问每个检查点；
  \item 每到一个检查点，机器人自动停留并通过相机采集现场画面；
  \item 完成全部检查点后返回出发位置，或继续下一轮巡检。
\end{enumerate}

\begin{infobox}[典型巡检场景]
以工厂消防检查为例：操作员在App地图上依次标记所有灭火器的位置作为检查点，设置每天巡检两次。灵途每天自动带着机器人走完全部检查点，通过相机拍摄每个灭火器的状态画面并回传给监控端。操作员坐在办公室里就能确认所有消防设施正常。
\end{infobox}


% ══════════════════════════════════════════════════════════════
%  第五章：语义导航
% ══════════════════════════════════════════════════════════════
\chapter{语义导航}

语义导航是灵途最具差异化的核心能力。它让用户可以用日常口语描述导航目标，而不是在地图上精确点击坐标。

\section{什么是语义导航}

传统导航系统的交互方式是这样的：用户在地图上找到目标位置 → 精确点击坐标 → 机器人前往。这要求用户必须知道目标的准确位置，操作门槛高。

灵途的语义导航改变了这一模式。用户只需用日常语言描述想要去哪里或想找什么，系统会自动理解语义、在环境中找到对应的物体或位置，并规划导航路线。

\begin{figure}[!htb]
\centering
\begin{tikzpicture}[
  cmd/.style={draw=qpblue!30, fill=cardblue, rounded corners=6pt, minimum width=6.2cm, minimum height=1.2cm, align=center, font=\small},
  result/.style={draw=qpgreen!30, fill=cardgreen, rounded corners=6pt, minimum width=6.2cm, minimum height=1.2cm, align=center, font=\small},
  arr/.style={-{Stealth[length=5pt]}, thick, qpblue!60},
]
  \node[font=\bfseries\color{qpdark}] at (-3, 3.5) {用户说：};
  \node[font=\bfseries\color{qpdark}] at (5.5, 3.5) {灵途执行：};

  \node[cmd] (c1) at (-1.5, 2.5) {"去大门那边看看"};
  \node[result] (r1) at (7, 2.5) {找到大门位置 → 自动导航前往};
  \draw[arr] (c1) -- (r1);

  \node[cmd] (c2) at (-1.5, 1) {"帮我找一下灭火器"};
  \node[result] (r2) at (7, 1) {识别最近的灭火器 → 导航到旁边};
  \draw[arr] (c2) -- (r2);

  \node[cmd] (c3) at (-1.5, -0.5) {"去机械设备那里检查"};
  \node[result] (r3) at (7, -0.5) {定位机械设备 → 前往并停留拍照};
  \draw[arr] (c3) -- (r3);

  \node[cmd] (c4) at (-1.5, -2) {"先去仓库，再去办公室"};
  \node[result] (r4) at (7, -2) {分解为两步任务 → 依次执行};
  \draw[arr] (c4) -- (r4);

  \node[cmd] (c5) at (-1.5, -3.5) {"跟着穿红衣服的人走"};
  \node[result] (r5) at (7, -3.5) {识别目标人物 → 持续跟随};
  \draw[arr] (c5) -- (r5);
\end{tikzpicture}
\caption{语义导航指令示例}
\end{figure}

\section{工作原理}

当用户发出一条语义导航指令后，系统经历四个阶段自动完成任务：

\begin{figure}[!htb]
\centering
\begin{tikzpicture}[
  stage/.style={draw=qpblue!40, fill=white, rounded corners=8pt, minimum width=13cm, minimum height=2.4cm, align=left, inner sep=10pt, font=\small},
  snum/.style={circle, fill=qpblue, text=white, font=\bfseries\large, minimum size=1cm},
]
  \node[stage] (s1) at (0, 7.5) {
    \hspace{1.5cm}{\bfseries\large\color{qpblue}指令理解} \\[3pt]
    \hspace{1.5cm}系统分析用户指令的含义，将复杂指令拆解为有序的子任务。\\
    \hspace{1.5cm}例如"先去大门，再找灭火器"被拆解为两个顺序执行的导航目标。\\
    \hspace{1.5cm}系统完全支持中文自然语言，内置了丰富的工业领域词汇库。
  };
  \node[snum] at ([xshift=0.7cm]s1.west) {1};

  \node[stage] (s2) at (0, 4.2) {
    \hspace{1.5cm}{\bfseries\large\color{qpblue}智能匹配} \\[3pt]
    \hspace{1.5cm}系统在已构建的环境认知模型中搜索与指令匹配的物体或位置。\\
    \hspace{1.5cm}超过75\%的指令可以在0.17秒内直接匹配成功（"快速通道"）。\\
    \hspace{1.5cm}对于模糊描述，系统自动启用AI深度推理（"深度通道"，约2秒）。
  };
  \node[snum] at ([xshift=0.7cm]s2.west) {2};

  \node[stage] (s3) at (0, 0.9) {
    \hspace{1.5cm}{\bfseries\large\color{qpblue}导航执行} \\[3pt]
    \hspace{1.5cm}确定目标坐标后，系统自动规划路线并执行导航。\\
    \hspace{1.5cm}行走过程中持续避障、检测卡死、监控进度。\\
    \hspace{1.5cm}到达目标附近后，系统会验证找到的确实是用户要找的东西。
  };
  \node[snum] at ([xshift=0.7cm]s3.west) {3};

  \node[stage] (s4) at (0, -2.4) {
    \hspace{1.5cm}{\bfseries\large\color{qpblue}智能恢复} \\[3pt]
    \hspace{1.5cm}如果导航中途遇到问题（路被堵、目标不对），系统自动三步恢复：\\
    \hspace{1.5cm}观察环境变化 → 分析失败原因 → 重新规划路线。\\
    \hspace{1.5cm}连续失败时，AI会决定：换条路、扩大搜索、重新理解指令、或向用户报告。
  };
  \node[snum] at ([xshift=0.7cm]s4.west) {4};

  \draw[-{Stealth[length=4pt]}, gray!40, thick] (s1.south) -- (s2.north);
  \draw[-{Stealth[length=4pt]}, gray!40, thick] (s2.south) -- (s3.north);
  \draw[-{Stealth[length=4pt]}, gray!40, thick] (s3.south) -- (s4.north);
\end{tikzpicture}
\caption{语义导航四阶段流程}
\end{figure}

\section{"快速+深度"双通道决策}

灵途独创的双通道决策架构，模仿了人类大脑的"直觉"与"深思"两种思维方式：

\begin{figure}[!htb]
\centering
\begin{tikzpicture}[
  chan/.style={rounded corners=10pt, minimum width=6cm, minimum height=4.5cm, align=center, inner sep=10pt},
]
  % 快速通道
  \node[chan, draw=qpblue!60, fill=qpblue!8] (fast) at (0, 0) {};
  \node[font=\Large\bfseries\color{qpblue}] at (0, 1.5) {快速通道};
  \node[font=\small, align=center] at (0, 0) {
    相当于人类的"一眼就看到了"\\[4pt]
    响应时间：\textbf{0.17 秒}\\[2pt]
    处理比例：\textbf{75\% 以上的指令}\\[2pt]
    无需网络，完全离线运行\\[4pt]
    {\color{qpgray}\footnotesize 适合简单、明确的指令}
  };

  % 深度通道
  \node[chan, draw=qporange!60, fill=qporange!8] (slow) at (8, 0) {};
  \node[font=\Large\bfseries\color{qporange}] at (8, 1.5) {深度通道};
  \node[font=\small, align=center] at (8, 0) {
    相当于人类的"仔细想一想"\\[4pt]
    响应时间：\textbf{约 2 秒}\\[2pt]
    处理比例：\textbf{25\% 的复杂指令}\\[2pt]
    调用云端AI进行深度推理\\[4pt]
    {\color{qpgray}\footnotesize 适合模糊、复杂的描述}
  };

  % 中间判断
  \node[diamond, draw=qpdark, fill=white, minimum size=1.2cm, font=\footnotesize, align=center] (dec) at (4, -3.5) {自动\\判断};
  \draw[-{Stealth[length=5pt]}, thick, qpblue] (dec) -- (0, -3.5) -- (fast.south);
  \draw[-{Stealth[length=5pt]}, thick, qporange] (dec) -- (8, -3.5) -- (slow.south);
  \node[font=\footnotesize\color{qpblue}, anchor=east] at (1.5, -3.1) {匹配度高};
  \node[font=\footnotesize\color{qporange}, anchor=west] at (6.5, -3.1) {匹配度低或模糊};

  % 输入
  \draw[-{Stealth[length=5pt]}, thick, qpdark] (4, -5) -- (dec);
  \node[draw=qpdark!30, fill=cardblue, rounded corners=4pt, font=\small] at (4, -5.5) {用户指令输入};
\end{tikzpicture}
\caption{"快速+深度"双通道决策架构}
\end{figure}

这种设计的好处是：简单指令几乎瞬间响应，不浪费时间；复杂指令也能正确理解，不会出错。系统会自动判断应该用哪个通道，用户无需关心。

\begin{tipbox}[即使没有网络也能工作]
快速通道完全在机器人本地运行，不依赖任何网络连接。只有深度通道在处理复杂语义时需要调用云端AI。这意味着在没有网络的工厂车间、地下室等环境中，灵途仍然可以处理大部分导航请求。
\end{tipbox}

\section{环境认知能力}

语义导航的基础，是机器人在行走过程中持续建立的\textbf{环境认知模型}（我们称之为"场景图"）。它不是一张简单的物体清单，而是一个有层次、有关系、有记忆的智能模型：

\begin{figure}[!htb]
\centering
\begin{tikzpicture}[
  level/.style={draw=qpblue!30, fill=#1, rounded corners=6pt, minimum width=11cm, minimum height=1.3cm, align=left, inner xsep=12pt, font=\small},
  arr/.style={-{Stealth[length=4pt]}, gray!40, thick},
]
  \node[level=cardblue!60] (l1) at (0, 4) {\textbf{楼层层}：代表一个物理楼层};
  \node[level=cardblue!40] (l2) at (0, 2.4) {\textbf{区域层}：自动识别区域类型（走廊、车间、仓库、办公室……）};
  \node[level=cardgreen!40] (l3) at (0, 0.8) {\textbf{功能组}：将相关物体聚合（桌+椅+电脑 = "办公工位"）};
  \node[level=cardorange!40] (l4) at (0, -0.8) {\textbf{物体层}：每个具体的物体，带有三维位置和可信度};

  \draw[arr] (l1.south) -- (l2.north);
  \draw[arr] (l2.south) -- (l3.north);
  \draw[arr] (l3.south) -- (l4.north);

  \node[font=\footnotesize\itshape\color{qpgray}, align=center] at (0, -2) {场景图在机器人行走时持续更新\\看到新物体就添加，重复看到就增强可信度，长时间不见就降低可信度};
\end{tikzpicture}
\caption{环境认知模型的层次结构}
\end{figure}

\section{三种记忆能力}

灵途的语义规划器具备三种互补的记忆能力，帮助机器人积累经验、做出更好的判断：

\begin{table}[!htb]
\centering
\renewcommand{\arraystretch}{1.6}
\begin{tabularx}{\textwidth}{>{\bfseries\color{qpdark}}l X >{\small\color{qpgray}}l}
  \toprule
  记忆类型 & 功能说明 & 容量 \\
  \midrule
  位置记忆 & 记录机器人到过的每一个关键位置，以及位置之间的连接关系。就像人类记住"从大门左转就是走廊"一样。 & 500个位置 \\
  经验记忆 & 以时间线记录在每个位置看到的物体。支持回忆"我之前在哪里见过灭火器"。 & 500条记录 \\
  常识库 & 预置了72个工业场景常见概念及其关联。如"灭火器通常在走廊墙上""安全出口标志在门上方"。即使没有直接看到目标，也能做出合理推测。 & 72个概念 \\
  \bottomrule
\end{tabularx}
\caption{灵途的三种记忆能力}
\end{table}

\section{自主探索}

在完全未知的环境中（没有预建地图），灵途的"探索模式"让机器人自主探索环境，同时寻找用户指定的目标。

系统不断识别已知区域与未知区域的边界，选择最有价值的方向前进。评估标准包括五个维度：距离远近、新颖程度、与目标的语言相关性、已识别物体暗示的区域类型、以及视觉多样性。

这就像一个经验丰富的搜救人员——他不会盲目乱跑，而是根据现有线索判断哪个方向最可能找到目标。

\section{人物跟随}

灵途支持让机器人跟随指定的人移动。操作非常自然：

\begin{enumerate}[leftmargin=2em, itemsep=2pt]
  \item 告诉机器人"跟着穿红衣服的人走"；
  \item 系统从画面中自动识别符合描述的人，记住他的外观特征；
  \item 机器人持续跟随目标人物，保持适当距离；
  \item 即使目标人物被短暂遮挡后重新出现，系统也能通过外观记忆重新锁定。
\end{enumerate}


% ══════════════════════════════════════════════════════════════
%  第六章：智能感知
% ══════════════════════════════════════════════════════════════
\chapter{智能感知}

感知系统是灵途的"眼睛"。它让机器人不仅能看到周围的环境，还能理解看到的是什么。

\section{传感器配置}

灵途搭载三种传感器，各自承担不同的感知任务：

\begin{figure}[!htb]
\centering
\begin{tikzpicture}[
  sensor/.style={draw=qpblue!30, fill=white, rounded corners=8pt, minimum width=4.2cm, minimum height=4.5cm, align=center, inner sep=8pt, font=\small},
]
  \node[sensor] (s1) at (0, 0) {
    {\bfseries\large\color{qpblue}激光雷达}\\[6pt]
    360°三维环境扫描\\[2pt]
    每秒20万个测量点\\[2pt]
    最远探测距离40米\\[2pt]
    用于建图和定位\\[6pt]
    {\color{qpgray}\footnotesize Livox Mid-360}
  };
  \node[sensor] (s2) at (5.5, 0) {
    {\bfseries\large\color{qpblue}深度相机}\\[6pt]
    同时拍摄彩色图像\\和距离信息\\[2pt]
    每秒30帧\\[2pt]
    用于物体识别\\[6pt]
    {\color{qpgray}\footnotesize Orbbec Gemini 335}
  };
  \node[sensor] (s3) at (11, 0) {
    {\bfseries\large\color{qpblue}姿态传感器}\\[6pt]
    测量机器人的\\姿态和角度\\[2pt]
    用于定位数据融合\\和倾斜安全保护\\[6pt]
    {\color{qpgray}\footnotesize 底盘内置}
  };
\end{tikzpicture}
\caption{灵途三种传感器}
\end{figure}

\section{物体识别能力}

灵途搭载了专门为工业巡检场景定制的AI视觉引擎，能够在机器人行走过程中实时识别环境中的物体。

\begin{figure}[!htb]
\centering
\begin{tikzpicture}[
  stp/.style={draw=qpblue!30, fill=cardblue, rounded corners=6pt, minimum width=2.2cm, minimum height=1.6cm, align=center, font=\footnotesize},
  arr/.style={-{Stealth[length=5pt]}, thick, qpblue!50},
]
  \node[stp] (s1) at (0, 0) {拍摄画面\\质量检查};
  \node[stp] (s2) at (3, 0) {AI物体\\识别};
  \node[stp] (s3) at (6, 0) {语义特征\\编码};
  \node[stp] (s4) at (9, 0) {三维位置\\计算};
  \node[stp] (s5) at (12, 0) {跨帧追踪\\去重合并};

  \draw[arr] (s1) -- (s2);
  \draw[arr] (s2) -- (s3);
  \draw[arr] (s3) -- (s4);
  \draw[arr] (s4) -- (s5);

  \node[font=\scriptsize\color{qpgray}, align=center] at (0, -1.3) {过滤因行走\\晃动导致的\\模糊画面};
  \node[font=\scriptsize\color{qpgray}, align=center] at (3, -1.3) {识别画面中\\所有物体\\含精确轮廓};
  \node[font=\scriptsize\color{qpgray}, align=center] at (6, -1.3) {计算每个物体\\的"语义指纹"\\用于语言匹配};
  \node[font=\scriptsize\color{qpgray}, align=center] at (9, -1.3) {将2D图像中\\的物体定位到\\3D地图坐标};
  \node[font=\scriptsize\color{qpgray}, align=center] at (12, -1.3) {同一物体在\\不同帧中只\\记录一次};

  \node[draw=qpgreen!50, fill=cardgreen, rounded corners=4pt, font=\small, minimum width=14cm] at (6, -3) {最终输出：一份持续更新的三维语义场景图，记录每个物体的名称、位置和可信度};
\end{tikzpicture}
\caption{物体识别五步管线}
\end{figure}

\subsection{识别范围}

系统预置了\textbf{125种}工业和户外场景常见物体，覆盖以下主要类别：

\begin{multicols}{3}
\begin{itemize}[leftmargin=1.5em, itemsep=1pt, parsep=0pt]
  \item 人员 / 行人
  \item 车辆 / 叉车
  \item 门 / 大门 / 安全门
  \item 灭火器 / 消火栓
  \item 安全出口标志
  \item 电气柜 / 配电箱
  \item 台阶 / 楼梯
  \item 栏杆 / 围栏
  \item 桌子 / 椅子
  \item 显示器 / 电脑
  \item 机械设备
  \item 管道 / 阀门
  \item 垃圾桶
  \item 植物 / 树木
  \item 交通标志 / 锥桶
  \item $\cdots$（共125类）
\end{itemize}
\end{multicols}

此外，系统还具备"开放词汇"能力——即使是未在训练集中出现过的物体名称，系统也能通过语义相似度进行模糊匹配。例如，用户说"找咖啡壶"，系统可能会匹配到最接近的"厨房设备"类物体。

\subsection{识别性能}

AI视觉引擎运行在专用AI加速芯片上（128 TOPS算力），实现了极高的实时性：

\begin{table}[!htb]
\centering
\renewcommand{\arraystretch}{1.4}
\begin{tabularx}{\textwidth}{X c c}
  \toprule
  \textbf{能力} & \textbf{处理速度} & \textbf{说明} \\
  \midrule
  通用物体识别（80类） & 每帧 22 毫秒 (45帧/秒) & 人、车、家具等常见物体 \\
  导航专用识别（125类） & 每帧 29 毫秒 (35帧/秒) & 含灭火器、电气柜等工业物体 \\
  场景图更新 & 每秒 1--2 次 & 持续更新环境认知模型 \\
  \bottomrule
\end{tabularx}
\caption{物体识别性能}
\end{table}

\section{地形感知}

除了识别物体，灵途还利用激光雷达数据实时分析地面地形，自动判断哪些区域可以安全通行：

\begin{itemize}[leftmargin=2em, itemsep=4pt]
  \item 高于地面 20 厘米以上的结构自动标记为障碍物；
  \item 地面坡度较大的区域增加通行难度权重，优先选择平坦路面；
  \item 地面突然下陷的位置（如台阶边缘、沟渠）被特别标记。
\end{itemize}

地形分析结果实时更新，直接影响机器人的路线选择——它会自动绕开障碍物和危险地形，选择最安全的行走路线。

\section{目标跟踪}

针对需要持续追踪运动目标的场景，灵途提供三种跟踪能力：

\begin{figure}[!htb]
\centering
\begin{tikzpicture}[
  track/.style={draw=qpblue!30, fill=#1, rounded corners=8pt, minimum width=4.5cm, minimum height=3.5cm, align=center, inner sep=8pt, font=\small},
]
  \node[track=cardblue] (t1) at (0, 0) {
    {\large\bfseries\color{qpblue} 实时跟踪}\\[6pt]
    处理速度：每秒23帧\\
    ID稳定性：100\%\\[4pt]
    适合通用巡检场景\\
    满足大多数需求
  };
  \node[track=cardorange] (t2) at (5.5, 0) {
    {\large\bfseries\color{qporange} 深度跟踪}\\[6pt]
    处理速度：每秒13帧\\
    含外观记忆功能\\[4pt]
    目标被遮挡后\\
    重新出现仍能识别
  };
  \node[track=cardpurple] (t3) at (11, 0) {
    {\large\bfseries\color{purple!70} 人物跟随}\\[6pt]
    语音选择跟随目标\\
    自动锁定并持续追踪\\[4pt]
    "跟着穿红衣服\\
    的那个人"
  };
\end{tikzpicture}
\caption{三种跟踪能力}
\end{figure}


% ══════════════════════════════════════════════════════════════
%  第七章：安全体系
% ══════════════════════════════════════════════════════════════
\chapter{安全体系}

安全是灵途系统的第一设计原则。我们不仅要让机器人"走得到"，更要确保它"走得安全"。灵途采用多层纵深防御策略，从硬件到软件共设置八道安全屏障。

\section{八层安全防护}

\begin{figure}[!htb]
\centering
\begin{tikzpicture}[
  lyr/.style={rounded corners=4pt, minimum width=13cm, minimum height=1.1cm, align=left, inner xsep=10pt, font=\small},
  pri/.style={circle, text=white, font=\bfseries\footnotesize, minimum size=0.7cm},
]
  \node[lyr, draw=qpred!60, fill=qpred!10] (l1) at (0, 7) {\hspace{1cm}\textbf{遥控器硬件覆盖} — 最高优先级，任何时刻一键接管};
  \node[pri, fill=qpred] at ([xshift=0.5cm]l1.west) {1};

  \node[lyr, draw=qpred!50, fill=qpred!7] (l2) at (0, 5.7) {\hspace{1cm}\textbf{控制板指令仲裁} — 遥控器活跃时，拒绝软件运动指令};
  \node[pri, fill=qpred!80] at ([xshift=0.5cm]l2.west) {2};

  \node[lyr, draw=qporange!60, fill=qporange!10] (l3) at (0, 4.4) {\hspace{1cm}\textbf{底盘看门狗} — 200毫秒内没收到新指令，自动停车};
  \node[pri, fill=qporange] at ([xshift=0.5cm]l3.west) {3};

  \node[lyr, draw=qporange!50, fill=qporange!7] (l4) at (0, 3.1) {\hspace{1cm}\textbf{安全门} — 速度上限1m/s、倾斜保护30度、近场避障};
  \node[pri, fill=qporange!80] at ([xshift=0.5cm]l4.west) {4};

  \node[lyr, draw=qpblue!50, fill=qpblue!7] (l5) at (0, 1.8) {\hspace{1cm}\textbf{模式状态机} — 非授权模式下禁止发布运动指令};
  \node[pri, fill=qpblue!80] at ([xshift=0.5cm]l5.west) {5};

  \node[lyr, draw=qpblue!40, fill=qpblue!5] (l6) at (0, 0.5) {\hspace{1cm}\textbf{健康监控} — 实时检查所有子系统，异常自动急停};
  \node[pri, fill=qpblue!60] at ([xshift=0.5cm]l6.west) {6};

  \node[lyr, draw=qpaccent!40, fill=qpaccent!5] (l7) at (0, -0.8) {\hspace{1cm}\textbf{地理围栏} — 设定工作区域边界，越界自动停车};
  \node[pri, fill=qpaccent!60] at ([xshift=0.5cm]l7.west) {7};

  \node[lyr, draw=qpgray!60, fill=qpgray!10] (l8) at (0, -2.1) {\hspace{1cm}\textbf{断连处理} — 与控制端失联30秒减速，5分钟停车};
  \node[pri, fill=qpgray] at ([xshift=0.5cm]l8.west) {8};

  % 标注
  \draw[decorate, decoration={brace, amplitude=8pt, mirror}, thick, qpred!60] (7, 7.5) -- (7, 5.2) node[midway, right=10pt, font=\footnotesize\color{qpred}] {硬件层保护};
  \draw[decorate, decoration={brace, amplitude=8pt, mirror}, thick, qporange!60] (7, 4.9) -- (7, 2.6) node[midway, right=10pt, font=\footnotesize\color{qporange}] {实时软件保护};
  \draw[decorate, decoration={brace, amplitude=8pt, mirror}, thick, qpblue!60] (7, 2.3) -- (7, 0) node[midway, right=10pt, font=\footnotesize\color{qpblue}] {系统级保护};
  \draw[decorate, decoration={brace, amplitude=8pt, mirror}, thick, qpaccent!60] (7, -0.3) -- (7, -2.6) node[midway, right=10pt, font=\footnotesize\color{qpaccent}] {运维级保护};
\end{tikzpicture}
\caption{八层纵深安全防护体系}
\end{figure}

\section{近场避障}

在遥控操作模式下，系统会实时分析前方地形，提供两级自动避障保护：

\begin{figure}[!htb]
\centering
\begin{tikzpicture}[
  zone/.style={minimum width=3cm, minimum height=2.5cm, rounded corners=4pt},
]
  % 机器人
  \node[draw=qpdark, fill=qpdark!20, minimum width=1.2cm, minimum height=1.5cm, rounded corners=3pt, font=\footnotesize\bfseries] (robot) at (0, 0) {机器人};

  % 红色急停区
  \node[zone, draw=qpred, fill=qpred!8, minimum width=2.5cm] (stop) at (3.5, 0) {};
  \node[font=\small\bfseries\color{qpred}] at (3.5, 0.7) {紧急制动区};
  \node[font=\footnotesize\color{qpred}] at (3.5, 0) {0 -- 0.8米};
  \node[font=\footnotesize\color{qpred}] at (3.5, -0.5) {立即停车};

  % 黄色减速区
  \node[zone, draw=qpwarn, fill=qpwarn!8, minimum width=3cm] (slow) at (7.5, 0) {};
  \node[font=\small\bfseries\color{qpwarn}] at (7.5, 0.7) {渐进减速区};
  \node[font=\footnotesize\color{qporange}] at (7.5, 0) {0.8 -- 2.0米};
  \node[font=\footnotesize\color{qporange}] at (7.5, -0.5) {按距离线性减速};

  % 绿色安全区
  \node[zone, draw=qpgreen, fill=qpgreen!8, minimum width=3cm] (safe) at (11.5, 0) {};
  \node[font=\small\bfseries\color{qpgreen}] at (11.5, 0.7) {安全区};
  \node[font=\footnotesize\color{qpgreen}] at (11.5, 0) {2.0米以外};
  \node[font=\footnotesize\color{qpgreen}] at (11.5, -0.5) {正常行驶};

  % 距离标注
  \draw[{Stealth[length=4pt]}-{Stealth[length=4pt]}, thick] (1, -2) -- (5, -2) node[midway, below, font=\footnotesize] {0.8m};
  \draw[{Stealth[length=4pt]}-{Stealth[length=4pt]}, thick] (5, -2) -- (10, -2) node[midway, below, font=\footnotesize] {1.2m};
\end{tikzpicture}
\caption{前方两级避障保护示意}
\end{figure}

即使操作员因视野盲区没有看到前方障碍物，机器人也会自动减速或停车。在紧急制动区内，线速度会立即归零，但角速度不受影响——机器人可以原地转向避让。

\section{三环实时监控}

灵途搭载了三个独立的实时监控环，像三位不同职责的安全员：

\begin{figure}[!htb]
\centering
\begin{tikzpicture}[
  rng/.style={rounded corners=10pt, minimum width=12cm, minimum height=2cm, align=left, inner xsep=12pt, font=\small},
]
  \node[rng, draw=qpred!60, fill=qpred!8] (r1) at (0, 4.5) {
    \textbf{\large\color{qpred} 安全监控}（每秒检查20次）\\[3pt]
    检查所有关键设备的心跳：激光雷达是否正常？底盘通信是否畅通？定位是否可信？\\
    发现异常立即按严重程度处理：轻微→预警，严重→减速，断裂→全停。响应时间50毫秒。
  };

  \node[rng, draw=qpblue!60, fill=qpblue!8] (r2) at (0, 1.5) {
    \textbf{\large\color{qpblue} 质量监控}（每秒检查5次）\\[3pt]
    评估导航执行效果：机器人是否偏离路线？前进速度是否正常？是否在原地打转？\\
    如果发现执行质量下降，自动通知规划系统重新规划路线。
  };

  \node[rng, draw=qpgreen!60, fill=qpgreen!8] (r3) at (0, -1.5) {
    \textbf{\large\color{qpgreen} 状态播报}（每秒更新2次）\\[3pt]
    将系统内部状态翻译成人类可读的信息，推送到用户App：\\
    "正在前往目标，已完成60\%，预计30秒后到达"、"前方发现障碍物，正在绕行"
  };
\end{tikzpicture}
\caption{三环实时监控体系}
\end{figure}

\begin{warnbox}[安全保证]
灵途的安全设计满足这样的极端假设：\textbf{即使导航软件完全崩溃、即使通信完全中断，机器人也一定会在200毫秒内自动停车}。遥控器作为最后的安全手段，通过硬件信号直连控制板，不经过任何软件处理——这是物理层面的安全保障。
\end{warnbox}


% ══════════════════════════════════════════════════════════════
%  第八章：远程管控
% ══════════════════════════════════════════════════════════════
\chapter{远程管控}

灵途为用户提供多种远程管控方式，让操作员坐在办公室就能完成机器人的所有操作。

\section{控制端App}

控制端App是灵途的主要操作界面，支持\textbf{Android手机}、\textbf{iOS}和\textbf{Windows电脑}三个平台。

\begin{figure}[!htb]
\centering
\begin{tikzpicture}[
  screen/.style={draw=qpdark!30, fill=white, rounded corners=3pt, minimum width=5.5cm, minimum height=7cm, inner sep=0pt},
  area/.style={draw=qpblue!20, fill=#1, rounded corners=2pt, minimum width=4.8cm, minimum height=1cm, align=center, font=\footnotesize},
  title/.style={font=\bfseries\small\color{qpdark}},
]
  % 手机框
  \node[screen] (phone) at (0, 0) {};
  \node[draw=qpdark!30, fill=qpdark!10, minimum width=5.5cm, minimum height=0.8cm, rounded corners=3pt, font=\footnotesize\bfseries\color{white}] at (0, 3.2) {\color{qpdark}灵途控制端};

  % App功能区域
  \node[area=cardblue] at (0, 2.2) {\title{实时地图} — 机器人位置 + 路径 + 目标};
  \node[area=cardgreen] at (0, 1) {\title{虚拟摇杆} — 遥控操作（含自动避障保护）};
  \node[area=cardorange] at (0, -0.2) {\title{语义导航} — 文字/语音指令输入};
  \node[area=cardpurple] at (0, -1.4) {\title{状态面板} — 速度/电量/姿态/健康};
  \node[area=cardteal] at (0, -2.6) {\title{模式切换} — 建图/导航/遥控/急停};

  % 右侧功能列表
  \node[font=\bfseries\color{qpdark}] at (7.5, 3) {完整功能列表};
  \node[font=\small, align=left, anchor=north west] at (5.5, 2.3) {
    \ding{51}~实时遥测（每秒30帧）\\[2pt]
    \ding{51}~三维地图显示与导航目标设定\\[2pt]
    \ding{51}~遥控操作（含安全保护）\\[2pt]
    \ding{51}~自然语言/语音导航指令\\[2pt]
    \ding{51}~多航点巡检任务配置\\[2pt]
    \ding{51}~运行模式切换\\[2pt]
    \ding{51}~地图保存/加载/管理\\[2pt]
    \ding{51}~系统健康诊断\\[2pt]
    \ding{51}~相机实时画面查看\\[2pt]
    \ding{51}~远程软件升级（OTA）\\[2pt]
    \ding{51}~多机器人管理\\[2pt]
    \ding{51}~操作日志记录
  };
\end{tikzpicture}
\caption{控制端App界面示意与功能列表}
\end{figure}

\begin{infobox}[双通道安全连接]
App同时通过两个独立通道连接机器人：一个连接导航系统，一个连接底盘控制系统。即使导航系统出现故障，用户仍然可以通过底盘直连通道进行遥控操作和传感器数据查看——\textbf{任何情况下都不会完全失去对机器人的控制}。
\end{infobox}

\section{Web监控仪表盘}

Web仪表盘是一个轻量级的浏览器端监控工具。无需安装任何软件，用同一网络内的任何设备打开浏览器访问即可使用。

主要用于运维人员快速查看机器人的整体运行状态，包括：各传感器数据更新频率、各服务进程运行状态、CPU/内存/存储使用率、以及语义导航的最近事件记录。

\section{远程软件升级}

灵途支持OTA（Over-The-Air）远程软件升级，让用户无需到现场就能升级机器人。

\begin{figure}[!htb]
\centering
\begin{tikzpicture}[
  ota/.style={draw=qpblue!30, fill=#1, rounded corners=6pt, minimum width=4.2cm, minimum height=2.5cm, align=center, font=\small, inner sep=6pt},
]
  \node[ota=cardgreen] (hot) at (0, 0) {
    {\large\bfseries\color{qpgreen} 热更新}\\[6pt]
    地图、配置文件\\[2pt]
    \textbf{无需停机}\\
    运行中即可安装
  };
  \node[ota=cardorange] (warm) at (5.5, 0) {
    {\large\bfseries\color{qporange} 温更新}\\[6pt]
    AI识别模型\\[2pt]
    \textbf{短暂暂停}\\
    约30秒后自动恢复
  };
  \node[ota=cardred] (cold) at (11, 0) {
    {\large\bfseries\color{qpred} 冷更新}\\[6pt]
    系统固件\\[2pt]
    \textbf{需进入维护模式}\\
    更新后自动重启
  };

  \node[font=\footnotesize\color{qpgray}, align=center] at (5.5, -2.2) {所有更新包经过数字签名验证，防止篡改\\安装过程支持断点续传和自动回滚，不会因断电导致系统损坏};
\end{tikzpicture}
\caption{三级OTA远程更新}
\end{figure}

\section{多机管理}

灵途支持同时管理多台机器人。每台机器人运行独立的灵途实例，通过不同的网络地址区分。控制端App支持在机器人列表中快速切换，一个操作员可以同时监控和管理多台机器人的运行状态和任务进度。


% ══════════════════════════════════════════════════════════════
%  第九章：系统架构
% ══════════════════════════════════════════════════════════════
\chapter{系统架构}

本章面向希望深入了解灵途技术架构的用户，介绍系统的整体设计和硬件部署方案。

\section{七层软件架构}

灵途采用七层架构设计。每一层专注于一个特定的功能领域，上层依赖下层提供的服务，但下层不依赖上层——这保证了系统的模块化和可靠性。

\begin{figure}[!htb]
\centering
\begin{tikzpicture}[
  layer/.style={draw=qpblue!30, fill=#1, rounded corners=3pt, minimum width=12cm, minimum height=1cm, align=center, font=\small},
]
  \node[layer=cardpurple!50] at (0, 6.5) {\textbf{交互层}：手机App + Web仪表盘 + 远程控制接口};
  \node[layer=cardred!30] at (0, 5.3) {\textbf{编排层}：任务管理 + 航点调度 + 巡逻/围栏等高级服务};
  \node[layer=cardorange!40] at (0, 4.1) {\textbf{规划层}：语义理解 + 全局路线规划 + 局部实时避障};
  \node[layer=cardgreen!40] at (0, 2.9) {\textbf{感知层}：AI物体识别 + 地形分析 + 目标跟踪};
  \node[layer=cardblue!40] at (0, 1.7) {\textbf{定位层}：激光定位建图 + 重定位};
  \node[layer=cardblue!20] at (0, 0.5) {\textbf{驱动层}：激光雷达 + 深度相机 + 底盘通信};
  \node[layer=white] at (0, -0.7) {\textbf{公共层}：数据校验 + 容错保护 + 全系统共享工具};
\end{tikzpicture}
\caption{灵途七层软件架构}
\end{figure}

\section{硬件部署}

灵途系统部署在机器人背部的计算板上，整体硬件拓扑如下：

\begin{figure}[!htb]
\centering
\begin{tikzpicture}[
  board/.style={draw=qpblue!40, fill=white, rounded corners=8pt, minimum width=7.5cm, minimum height=4cm, align=center, inner sep=8pt},
  btitle/.style={font=\bfseries\large\color{qpblue}},
  item/.style={font=\small, align=left},
  client/.style={draw=qpaccent!40, fill=cardteal, rounded corners=8pt, minimum width=5cm, minimum height=2cm, align=center, font=\small},
  conn/.style={<->, thick, qpblue!60},
]
  % 导航板
  \node[board] (nav) at (0, 0) {};
  \node[btitle] at (0, 1.5) {导航计算板};
  \node[item] at (0, 0) {
    AI加速芯片 128 TOPS\\
    16GB内存 + 238GB存储\\
    激光雷达 + 深度相机\\
    运行灵途全部导航软件
  };

  % 控制板
  \node[board, minimum height=3cm] (dog) at (9, 0) {};
  \node[btitle] at (9, 1) {运动控制板};
  \node[item] at (9, -0.3) {
    运动控制策略\\
    16关节电机控制\\
    姿态传感器
  };

  % 连接
  \draw[conn] (nav.east) -- (dog.west) node[midway, above, font=\footnotesize\color{qpgray}] {高速通信};

  % 客户端
  \node[client] (app) at (0, -4.5) {
    {\bfseries\color{qpaccent} 控制端App}\\
    手机 / 电脑
  };
  \node[client] (rc) at (9, -4.5) {
    {\bfseries\color{qpaccent} 遥控器}\\
    硬件直连（最高优先级）
  };

  \draw[conn] (nav.south) -- (app.north) node[midway, right, font=\footnotesize\color{qpgray}] {WiFi};
  \draw[conn] (dog.south) -- (rc.north) node[midway, right, font=\footnotesize\color{qpgray}] {信号线};

  % 传感器标注
  \node[draw=qpgray!30, fill=scenebg, rounded corners=3pt, font=\footnotesize, inner sep=4pt, align=center] at (-4.5, 0) {激光雷达\\深度相机\\姿态传感器};
  \draw[-{Stealth[length=4pt]}, qpgray!60] (-3.2, 0) -- (nav.west);
\end{tikzpicture}
\caption{灵途硬件部署拓扑}
\end{figure}

导航板和控制板的分离设计确保了：导航软件升级不影响运动控制的稳定性；即使导航板出现故障，遥控器仍然可以直接控制机器人。

\section{运行模式}

灵途提供三种运行模式，覆盖不同的使用阶段：

\begin{table}[!htb]
\centering
\renewcommand{\arraystretch}{1.6}
\begin{tabularx}{\textwidth}{>{\bfseries\color{qpdark}}l c c c X}
  \toprule
  模式 & 预建地图 & 自主导航 & 语义理解 & 适用场景 \\
  \midrule
  建图模式 & 不需要 & 否 & 否 & 首次到达新环境，遥控走一圈构建地图 \\
  导航模式 & 需要 & 是 & 可选 & 日常工作，在已知环境中自主/遥控导航 \\
  探索模式 & 不需要 & 是 & 是 & 未知环境自主探索，边走边建图边找目标 \\
  \bottomrule
\end{tabularx}
\caption{三种运行模式}
\end{table}


% ══════════════════════════════════════════════════════════════
%  第十章：技术规格
% ══════════════════════════════════════════════════════════════
\chapter{技术规格}

\section{性能参数}

\begin{table}[!htb]
\centering
\renewcommand{\arraystretch}{1.4}
\begin{tabularx}{\textwidth}{X c c}
  \toprule
  \textbf{指标} & \textbf{设计目标} & \textbf{实测值} \\
  \midrule
  \multicolumn{3}{l}{\textbf{定位与建图}} \\
  \quad 定位更新频率 & $\geq$ 10 次/秒 & 10 次/秒 \\
  \quad 建图精度（标准模式） & 误差 $<$ 5 厘米 & 0.8 厘米 \\
  \quad 建图精度（增强抗振模式） & 误差 $<$ 5 厘米 & 0.6 厘米 \\
  \midrule
  \multicolumn{3}{l}{\textbf{物体识别}} \\
  \quad 通用场景识别速度 & $<$ 50 毫秒/帧 & 22 毫秒 (45帧/秒) \\
  \quad 导航专用识别速度 & $<$ 50 毫秒/帧 & 29 毫秒 (35帧/秒) \\
  \quad 识别类别数 & $\geq$ 80 类 & 125 类 \\
  \midrule
  \multicolumn{3}{l}{\textbf{路径规划}} \\
  \quad 路线规划速度 & $<$ 50 毫秒 & 5--10 毫秒 \\
  \midrule
  \multicolumn{3}{l}{\textbf{语义导航}} \\
  \quad 快速通道响应时间 & $<$ 200 毫秒 & 约 170 毫秒 \\
  \quad 快速通道命中率 & $>$ 70\% & 约 75\% \\
  \quad 深度通道响应时间 & $<$ 5 秒 & 约 2 秒 \\
  \midrule
  \multicolumn{3}{l}{\textbf{目标跟踪}} \\
  \quad 实时跟踪帧率 & $>$ 15 帧/秒 & 23 帧/秒 \\
  \quad 深度跟踪帧率 & $>$ 10 帧/秒 & 13 帧/秒 \\
  \quad 跟踪 ID 稳定性 & $>$ 95\% & 100\% \\
  \midrule
  \multicolumn{3}{l}{\textbf{系统}} \\
  \quad 开机到就绪时间 & $<$ 60 秒 & 约 30 秒 \\
  \quad 全链路端到端延迟 & $<$ 60 秒 & 18 秒 \\
  \bottomrule
\end{tabularx}
\caption{灵途系统性能参数}
\end{table}

\section{机器人参数}

\begin{table}[!htb]
\centering
\renewcommand{\arraystretch}{1.4}
\begin{tabularx}{\textwidth}{>{\bfseries}l l X}
  \toprule
  参数 & 数值 & 说明 \\
  \midrule
  车身尺寸 & 100 × 60 × 50 cm & 长 × 宽 × 高 \\
  最大行走速度 & 1.0 m/s & 软件限速，可根据场景调整 \\
  自主巡航速度 & 0.875 m/s & 默认速度，兼顾效率与安全 \\
  紧急制动距离 & 0.8 米 & 前方障碍触发立即停车 \\
  渐进减速距离 & 2.0 米 & 前方障碍触发自动减速 \\
  底盘看门狗超时 & 0.2 秒 & 通信中断后自动停车时间 \\
  关节数量 & 16 个 & 四条腿，每条腿4个关节 \\
  \bottomrule
\end{tabularx}
\caption{机器人物理参数}
\end{table}

\section{计算平台}

\begin{table}[!htb]
\centering
\renewcommand{\arraystretch}{1.4}
\begin{tabularx}{\textwidth}{>{\bfseries}l X}
  \toprule
  参数 & 规格 \\
  \midrule
  AI算力 & 128 TOPS（专用AI加速芯片） \\
  内存 & 16 GB \\
  存储 & 238 GB 高速固态硬盘 \\
  网络 & WiFi + 以太网 \\
  操作系统 & Ubuntu 22.04 \\
  \bottomrule
\end{tabularx}
\caption{计算平台规格}
\end{table}


% ══════════════════════════════════════════════════════════════
%  第十一章：常见问题
% ══════════════════════════════════════════════════════════════
\chapter{常见问题}

\begin{enumerate}[leftmargin=2em, itemsep=14pt]

\item \textbf{灵途需要联网才能工作吗？}

不需要。灵途的所有核心功能（建图、导航、避障、物体识别）都在机器人本地运行，不依赖任何网络。仅在使用语义导航的"深度通道"处理复杂自然语言指令时需要网络来调用云端AI。即使没有网络，"快速通道"仍可处理 75\% 以上的导航请求。

\item \textbf{建一次图需要多长时间？}

取决于区域面积。一个 2000 平方米的标准工厂车间，操控机器人以正常步速行走一圈，通常 15--20 分钟即可完成。地图构建后保存在本地，后续使用无需重建。

\item \textbf{能在户外使用吗？}

可以。灵途专为户外和半户外环境设计，可以应对草地、碎石路、缓坡等常见户外地形。但不建议在大雨、积水或极端泥泞环境中使用。

\item \textbf{机器人失控了怎么办？}

使用遥控器——它通过硬件信号直连控制板，优先级最高，在任何情况下都能立即接管。此外，即使完全断开通信，底盘看门狗会在 0.2 秒内自动停车。灵途共有八层安全防护，确保任何单点故障都不会导致失控。

\item \textbf{能同时管理多台机器人吗？}

可以。每台机器人运行独立的灵途实例，通过不同的网络地址区分。控制端App支持在机器人列表中切换，一个操作员可以同时监控管理多台机器人。

\item \textbf{地图需要定期更新吗？}

如果环境发生了较大结构性变化（如搬迁大型设备、修建新墙），建议重新建图。日常小变化（如移动椅子），系统的实时避障功能可以自动应对。

\item \textbf{能识别多少种物体？}

系统预置了 125 种工业/户外场景常见物体。同时支持"开放词汇"能力——即使是未预置的物体名称，系统也能通过语义相似度进行模糊匹配。

\item \textbf{如何升级软件？}

灵途支持远程OTA升级。通过App检查更新并一键安装，无需到现场。地图和配置可以在运行中热更新；AI模型更新需短暂暂停（约30秒）；固件更新需进入维护模式。

\item \textbf{安装灵途需要对机器人做改装吗？}

灵途以独立计算模块的形式安装在机器人背部，只需进行传感器安装和线缆连接，不需要对机器人本体做任何结构改装。

\item \textbf{支持哪些品牌的四足机器人？}

灵途目前支持穹沛科技自研的四足机器人平台。如需适配其他品牌的四足机器人，请联系我们的技术团队评估。

\end{enumerate}


% ──────────────────────────────────────────────────────────────
% 封底
% ──────────────────────────────────────────────────────────────
\clearpage
\thispagestyle{empty}
\vspace*{\fill}
\begin{center}
  {\Large\color{qpblue}\textbf{上海穹沛科技有限公司}}

  \vspace{0.5cm}

  {\small\color{qpgray}INOVXIO Technology Co., Ltd.}

  \vspace{1.5cm}

  {\large\color{qpdark}灵途自主导航系统}

  \vspace{0.3cm}

  {\color{qpgray}让四足机器人自主思考、自主行动}

  \vspace{2.5cm}

  {\color{qpgray}
  技术咨询与商务合作 \\[0.5em]
  邮箱：contact@inovxio.com \\[0.3em]
  地址：上海市
  }

  \vspace{2cm}

  {\small\color{qpgray!60}
  \copyright\ 2026 上海穹沛科技有限公司 版权所有 \\
  本文档内容未经授权不得复制或传播
  }
\end{center}
\vspace*{\fill}

\end{document}
